% Generated draft: requires mathematical and source review.
\documentclass[11pt]{article}
\usepackage[margin=25mm]{geometry}
\usepackage{fontspec}
\setmainfont{Latin Modern Roman}
\usepackage{amsmath,amssymb,mathtools,unicode-math}
\setmathfont{Latin Modern Math}
\usepackage{enumitem,needspace,xurl,hyperref}
\setlength{\parindent}{0pt}
\setlength{\parskip}{5pt}
\setlength{\emergencystretch}{3em}
\newcommand{\N}{\mathbb{N}}
\newcommand{\Z}{\mathbb{Z}}
\newcommand{\Q}{\mathbb{Q}}
\newcommand{\R}{\mathbb{R}}
\newcommand{\C}{\mathbb{C}}
\newcommand{\F}{\mathbb{F}}
\newcommand{\A}{\mathbb{A}}
\newcommand{\PP}{\mathbb{P}}
\newcommand{\E}{\mathbb{E}}
\newcommand{\eps}{\varepsilon}
\newcommand{\dd}{\,\mathrm{d}}
\begin{document}
\section*{UnsolvedMath ID: 30002661}
\textbf{Source problem code:} catalogue entry
\subsection{Definitions and mathematical statement}
Let \(n,m\geq 1\) be integers. A convex body \(K\subset\R^n\) means a compact convex set with nonempty interior; no symmetry is assumed. Define its difference body and the associated norm by
\[
D=K-K=\{x-y:x,y\in K\},\qquad
\|v\|_D=\inf\{t>0:v\in tD\}.
\]
For a nonzero linear functional \(\ell:\R^n\to\R\), put
\[
w_K(\ell)=\max_{x\in K}\ell(x)-\min_{x\in K}\ell(x)>0.
\]
A plank is a closed region
\[
P=\{x\in\R^n:a\leq\ell(x)\leq b\},\qquad a<b.
\]
Interpret its relative width as the distance between its boundary hyperplanes in the norm with unit ball \(D\). Explicitly,
\[
r_K(P)=\inf_{\ell(v)=b-a}\|v\|_D
=\frac{b-a}{\max_{z\in D}|\ell(z)|}
=\frac{b-a}{w_K(\ell)}.
\]
The supplied record does not explicitly define relative width; this norm-distance interpretation fixes the normalization and requires confirmation against the cited report.

The question, identified in the supplied record as Bang's plank conjecture for nonsymmetric convex bodies, is whether for every such \(K\) and every finite family of planks \(P_1,\ldots,P_m\),
\[
K\subseteq\bigcup_{i=1}^{m}P_i
\quad\Longrightarrow\quad
\sum_{i=1}^{m}r_K(P_i)\geq 1.
\]
In particular, the question includes bodies that are not centrally symmetric. The supplied open-status label is a dated curation claim, not an independently verified mathematical conclusion.
\subsection{Short English statement}
Must every finite plank covering of a convex body have total relative width at least \(1\), without assuming central symmetry? Here each relative width is the distance between the plank's boundary hyperplanes in the norm whose unit ball is \(K-K\); this interpretation of the source terminology requires review.
\subsection{Sources}
\begin{itemize}
\item Local UnsolvedMath registry, \texttt{data/unsolvedmath/problems.json}, permanent ID \(30002661\), source problem code \texttt{catalogue entry}, title \emph{Bang Plank Conjecture for Nonsymmetric Convex Bodies}. The supplied registry records proposed year 2014 and gives no individual proposer.
\item Supplied original-report citation: \emph{Discrete Geometry (2015)}, Oberwolfach Reports, pp.\ 2235--2296, DOI \texttt{10.4171/owr/2014/40}: \url{https://doi.org/10.4171/owr/2014/40}. The supplied background attributes the problem collection to workshop open-problem discussions documented in Oberwolfach Reports.
\item Supplied literature reference: Akopyan, Karasev, and Petrov, described in the record as ``symplectic invariants'': \url{https://doi.org/10.4310/JSG.2019.v17.n6.a1}. A complete title and a specific supporting location are not supplied.
\item Literature assessment dated exactly \texttt{2026-08-21}: supplied dated curation, not independently verified. It labels the problem open, states that the nonsymmetric-body version remains open, and reports that no later general resolution was found. The record lists the two references above as checked and explicitly requests independent verification of time-sensitive status and cited claims.
\end{itemize}
\end{document}
